\chapter{Literature Survey}
\label{ch:lit}

\section{Literature Survey}
This chapter summarises the research that Verif-AI builds on and the open problems that shaped its evaluation. It is not a complete survey;
broader reviews are given by Verdoliva \citep{verdoliva2020media} and Tolosana et al.\ \citep{tolosana2020deepfakes}.

\subsection{Evolution of Deepfake Generation and Detection}
\textbf{Generation.} Generative adversarial networks \citep{goodfellow2014gan} made photorealistic face synthesis practical;
StyleGAN and its successors generate high-resolution faces from a latent code \citep{karras2019stylegan,karras2020stylegan2}. Face-swapping
and re-enactment methods transfer one person's face or expression onto another's video, and the FaceForensics++ benchmark collects four
such manipulation families \citep{rossler2019ff}. Diffusion models \citep{ho2020ddpm}, in particular latent diffusion
\citep{rombach2022ldm}, now produce images from text prompts and can inpaint or edit regions of real photographs, so a single dataset
may contain fully synthetic, edited and swapped faces.

\textbf{Detection.} Early detectors used compact convolutional networks on face crops, such as MesoNet \citep{afchar2018mesonet}, or
general-purpose backbones such as Xception \citep{chollet2017xception} that performed best in the FaceForensics++ study
\citep{rossler2019ff}. Two-stream networks combine an RGB stream with a stream that sees low-level traces: Zhou et al.\
\citep{zhou2017twostream} paired a face-classification stream with a steganalysis-features stream for tampered faces, and constrained
convolutional layers \citep{bayar2016constrained} learn high-pass filters for manipulation detection. Dang et al.\ \citep{dang2020detection}
add attention maps that localise manipulated regions. Large benchmark efforts such as Celeb-DF \citep{li2020celebdf} and the DFDC
collection \citep{dolhansky2020dfdc} made video-level evaluation possible.

\subsection{Theoretical Framework of Modern Forensics}
\textbf{RGB versus noise residuals.} A generator leaves statistical traces in the high-frequency content of an image that are largely
independent of the depicted face. Rich-model steganalysis \citep{fridrich2012srm} extracts such traces by applying fixed high-pass
\emph{residual} filters and analysing the co-occurrence of the residuals. Verif-AI uses a fixed bank of three $3\times3$ high-pass
kernels (Chapter~\ref{ch:design}) that suppress image content and leave a noise residual for a second convolutional stream. The filters
are not trained, which keeps the stream's input independent of the training data.

\textbf{EfficientNet.} EfficientNet \citep{tan2019efficientnet} scales depth, width and resolution of a convolutional network with a
compound coefficient. The smallest member, EfficientNet-B0, has about 5.3 million parameters and reports 77.1\% ImageNet top-1
accuracy, which makes two parallel copies affordable on a CPU.

\textbf{Grad-CAM.} Gradient-weighted Class Activation Mapping \citep{selvaraju2017gradcam} weights the activation maps of a chosen
convolutional layer by the average gradient of the class score with respect to them and shows the result as a coarse heat-map. The map
indicates regions that increased the score; it is an explanation of the \emph{model}, not evidence that those pixels were altered
\citep{selvaraju2017gradcam}. Figure~\ref{fig:srm-concept} contrasts the information seen by the two streams.

\begin{figure}[ht]
\centering
\begin{adjustbox}{max width=\linewidth,center}
\begin{tikzpicture}[font=\small,x=1cm,y=1cm]
\foreach \i/\t/\c in {0/{RGB face (what a person sees)}/navy!10, 1/{High-pass residual (what the SRM stream sees)}/accent!10}{
  \node[draw,minimum width=3.7cm,minimum height=3.0cm,fill=\c] at (4.4*\i,0) {};
  \node[align=center,font=\footnotesize,text width=3.6cm] at (4.4*\i,-2.1) {\t};
}
\foreach \a in {0,...,7}{\foreach \b in {0,...,5}{
  \pgfmathtruncatemacro{\w}{20+mod(\a*7+\b*13,5)*10}
  \fill[navy!\w] ({-1.8+\a*0.46},{1.35-\b*0.5}) rectangle ++(0.42,0.46);
}}
\foreach \a in {0,...,31}{\foreach \b in {0,...,23}{
  \pgfmathtruncatemacro{\v}{(mod(\a*17+\b*29+\a*\b,7)>5) ? 40 : 6}
  \fill[accent!\v] ({2.6+\a*0.115},{1.4-\b*0.12}) rectangle ++(0.1,0.1);
}}
\draw[-{Stealth},thick] (1.9,0)--(2.5,0);
\end{tikzpicture}
\end{adjustbox}
\caption[Illustration (schematic, not a real face) of the idea behind the two streams: the RGB stream]{Illustration (schematic, not a real face) of the idea behind the two streams: the RGB stream sees image content, whereas the fixed
SRM-style filters remove content and leave a high-frequency noise residual in which generator traces may remain.}
\label{fig:srm-concept}
\end{figure}

\subsection{Analysis of Benchmark Architectures}
Table~\ref{tab:backbones} compares backbones commonly used for face-manipulation detection. The figures are the sizes and ImageNet
results reported by the original papers; they describe the architectures, not their forensic accuracy, which depends on training data.

\begin{table}[ht]
\centering
\caption{Backbones used in the deepfake-detection literature (published figures).}
\label{tab:backbones}
\small
\begin{adjustbox}{max width=\linewidth,center}
\begin{tabular}{|l|c|c|p{5.2cm}|}
\hline
\textbf{Architecture} & \textbf{Params (M)} & \textbf{ImageNet top-1} & \textbf{Relevance} \\ \hline
MesoNet (Meso-4) \citep{afchar2018mesonet} & 0.028 & -- & Very small network designed for face-video forgery \\ \hline
Xception \citep{chollet2017xception} & 22.9 & 79.0\% & Strongest baseline in FaceForensics++ \citep{rossler2019ff} \\ \hline
ResNet-50 \citep{he2016resnet} & 25.6 & -- & Common general-purpose baseline \\ \hline
\textbf{EfficientNet-B0} \citep{tan2019efficientnet} & \textbf{5.3} & \textbf{77.1\%} & Used for both streams of Verif-AI \\ \hline
\end{tabular}
\end{adjustbox}
\end{table}

\subsection{Comparison of Forensic Datasets}
Table~\ref{tab:datasets} lists the datasets that matter for this project. Their \emph{construction} matters as much as their size: how real and
fake samples are paired, whether identities repeat across splits, and how fakes were compressed or processed all create shortcuts that a
detector can learn instead of the manipulation itself (Section~\ref{sec:shortcuts}).

\begin{table}[ht]
\centering
\caption[Datasets referenced in this work (figures as stated by the dataset papers or by the dataset]{Datasets referenced in this work (figures as stated by the dataset papers or by the dataset documentation in the repository).}
\label{tab:datasets}
\small
\begin{adjustbox}{max width=\linewidth,center}
\begin{tabular}{|p{2.6cm}|p{3.0cm}|p{2.9cm}|p{3.4cm}|}
\hline
\textbf{Dataset} & \textbf{Content} & \textbf{Fakes} & \textbf{Use in Verif-AI} \\ \hline
FaceForensics++ \citep{rossler2019ff} & 1{,}000 original videos & 4 manipulation methods & Cited; not used (access agreement) \\ \hline
DFDC \citep{dolhansky2020dfdc} & over 100{,}000 clips & Several methods & Cited; not used \\ \hline
Celeb-DF v2 \citep{li2020celebdf} & 590 real and 5{,}639 fake videos & Face swap (one method) & Video-level training and test (research use only) \\ \hline
DeepFakeFace \citep{song2023dff} & 30{,}000 real (IMDB-WIKI) photographs \citep{rothe2018imdbwiki} & 3$\times$30{,}000: inpainting, face swap, text-to-image & Primary training data (sampled subset) \\ \hline
FFHQ \citep{karras2019stylegan} & 70{,}000 real faces & -- & Real faces in experiment E3 \\ \hline
SFHQ \citep{sfhq2022} & Synthetic faces & StyleGAN-generated & Fake faces in experiment E3 \\ \hline
\end{tabular}
\end{adjustbox}
\end{table}

\subsection{Generalisation, Shortcuts and Calibration}
\label{sec:shortcuts}
\textbf{Generalisation.} A classifier trained on images from one generator can be nearly perfect on that generator and weak on others.
Wang et al.\ \citep{wang2020cnn} found that a classifier trained on one GAN with suitable augmentation generalised surprisingly well to
other CNN-based generators, while stressing the limits of that result; Ojha et al.\ \citep{ojha2023universal} and Corvi et al.\ \citep{corvi2023diffusion} study the transfer between GAN and diffusion generators, where detectors trained on one family often degrade on
the other. A reliable evaluation therefore needs a held-out \emph{generator}, not only held-out images.

\textbf{Shortcut learning.} Geirhos et al.\ \citep{geirhos2020shortcut} describe how networks solve benchmarks with unintended cues.
In deepfake data the cue may be image sharpness, resolution, file format, compression history or the source of the real class.
Chapter~\ref{ch:results} documents a concrete case in this project.

\textbf{Calibration and explanation.} The sigmoid output of a modern network is usually over-confident
\citep{guo2017calibration}; Verif-AI therefore labels it a \emph{model score} and never a probability. Grad-CAM maps are explanations of
the model's decision, and as noted above they do not prove manipulation.

\subsection{Synthesis and Research Gap}
Prior work supplies the ingredients --- an efficient backbone, a noise-residual stream, class-activation explanations and benchmark
datasets --- but a practical forensic tool also needs a consistent preprocessing pipeline, video aggregation, privacy-preserving
handling of uploads and an evaluation that distinguishes internal from cross-generator behaviour. The gap addressed in this report is
therefore less a new architecture than a \emph{complete and honestly evaluated system}: the two-stream design is kept, while the data
splits, leakage checks, failure diagnosis and the limits of each claim are treated as first-class results.


\section{Research Gap, Project Contributions, and Novelty}
\label{sec:novelty}
This section states what existing work has achieved, which gaps matter for this project, which components are borrowed, and what Verif-AI itself contributes. \textbf{Verif-AI does not propose a new deep-learning architecture}, and it does not claim to have invented EfficientNet, SRM filtering, MTCNN or Grad-CAM; its contributions are in engineering, system integration and experimental method.

\subsection{What existing research has achieved}
Benchmarks and generators: FaceForensics++ \citep{rossler2019ff}, Celeb-DF \citep{li2020celebdf}, the DFDC dataset \citep{dolhansky2020dfdc} and DeepFakeFace \citep{song2023dff} supply labelled data. Detectors: compact and Xception-style CNNs \citep{afchar2018mesonet,rossler2019ff}; two-stream networks that combine an RGB branch with a low-level residual branch \citep{zhou2017twostream}; detectors that exploit high-frequency/SRM-filtered features to generalise across forgeries \citep{luo2021highfreq}; blending-boundary cues for unseen manipulations \citep{li2020facexray}; recurrent video models over frame features \citep{guera2018rnn,masi2020twobranch}; and detectors studied across generators \citep{wang2020cnn,ojha2023universal,corvi2023diffusion}. Explanation: Grad-CAM \citep{selvaraju2017gradcam}. Evaluation caveats: shortcut learning \citep{geirhos2020shortcut} and over-confident scores \citep{guo2017calibration}. Surveys: \citep{verdoliva2020media,tolosana2020deepfakes}.

\subsection{Gaps that matter for this project}
\begin{enumerate}[label=G\arabic*.]
\item The cited papers present \emph{methods and benchmark results}; they are not described as ready-to-use tools that return per-face evidence, video timelines and a record of which model and file produced a result.
\item Reported accuracy depends on the split. Identity-disjoint, duplicate-free and generator-disjoint testing, and the check that an apparent skill is not a quality shortcut, must be done again for any new system.
\item Video support is usually a temporal classifier validated inside one dataset; a deployed tool also needs face tracking, aggregation and honest handling of frames for which no evidence was stored.
\item A prototype inherited by this project (Forensic-Net, Section~\ref{sec:prototype}) reported very high accuracy that could not be reproduced on independent data, so the model itself needed to be re-trained and re-evaluated.
\end{enumerate}

\subsection{Established components used (not new)}
EfficientNet-B0 \citep{tan2019efficientnet} as both backbones; the three SRM-style high-pass kernels after Fridrich and Kodovsk\'y \citep{fridrich2012srm}; the two-stream idea of an image branch plus a low-level residual branch \citep{zhou2017twostream}; MTCNN \citep{zhang2016mtcnn}; Grad-CAM \citep{selvaraju2017gradcam}; ImageNet initialisation, AdamW \citep{loshchilov2019adamw} and PyTorch \citep{paszke2019pytorch}; and the public datasets of Chapter~\ref{ch:results}.

\subsection{The prototype and what Verif-AI adds}
\label{sec:prototype}
The RGB + three-kernel SRM two-stream design, with its feature concatenation and sigmoid head, comes from the earlier Forensic-Net prototype written by Lakshya Singh (notebook, checkpoint and plots preserved in the repository). That prototype analysed only the first frame of a file and had no application, no leakage-controlled evaluation and no tests. Everything listed below was developed for Verif-AI.

\begin{table}[ht]
\centering
\caption[Existing work and how Verif-AI relates to it (only capabilities stated in the cited papers or]{Existing work and how Verif-AI relates to it (only capabilities stated in the cited papers or in this repository).}
\label{tab:novelty}
\small
\begin{tabular}{|p{3.3cm}|p{3.9cm}|p{5.6cm}|}
\hline
\textbf{Existing work} & \textbf{What it contributes} & \textbf{Verif-AI: implemented / extended / different} \\ \hline
Zhou et al.\ \citep{zhou2017twostream} & two-stream face-tampering detector (image + low-level branch) & uses the two-stream idea with fixed SRM residuals; adds identity-disjoint training/evaluation and a per-quality-bin gate \\ \hline
Fridrich \& Kodovsk\'y \citep{fridrich2012srm} & rich residual models for steganalysis & three fixed kernels as the noise stream; examines when that stream carries a quality shortcut (Section~\ref{sec:diag}) \\ \hline
Luo et al.\ \citep{luo2021highfreq}; Li et al.\ \citep{li2020facexray} & generalisation through high-frequency features / blending boundaries & not re-implemented; the project measures generalisation with an unseen-generator split instead \\ \hline
G\"uera \& Delp \citep{guera2018rnn}; Masi et al.\ \citep{masi2020twobranch} & recurrent networks over frame features for video & different approach: per-face scores, IoU face tracks, per-track median, suspicious segments and stored evidence frames (unvalidated aggregation rule) \\ \hline
Selvaraju et al.\ \citep{selvaraju2017gradcam} & Grad-CAM & applied to the RGB stream, shown with Original/Heatmap/Overlay views, stored for evidence frames, stated as not proof \\ \hline
Wang et al., Ojha et al., Corvi et al.\ \citep{wang2020cnn,ojha2023universal,corvi2023diffusion} & detectors studied across generators & motivates the internal vs.\ cross-generator reporting; no claim of universal detection \\ \hline
Celeb-DF, DeepFakeFace, FFHQ/SFHQ \citep{li2020celebdf,song2023dff,karras2019stylegan,sfhq2022} & datasets & used under their licences; leakage control, near-duplicate removal and a licence review before use \\ \hline
\end{tabular}
\end{table}

\subsection{Contributions, by type}
\begin{description}
\item[Research novelty.] None claimed at the level of a new architecture or algorithm. The research-level results are empirical findings: (i) the production model was at chance on held-out Celeb-DF video; (ii) image sharpness is a measurable shortcut in DeepFakeFace and the noise stream carries it; (iii) a per-quality-bin acceptance gate (validated with synthetic controls) fails all four candidate models although two pass the usual ROC-AUC screen.
\item[Engineering novelty.] Training-matched preprocessing (JPEG round-trip), strict checkpoint loading with hash recording, an INCONCLUSIVE band and decisiveness value, atomic job storage, privacy-preserving deletion, memory-aware lite mode, and a model registry (default / alternative / rollback) with per-result model identification.
\item[System-level integration.] One React + FastAPI application that turns the two-stream model into an inspectable workflow: face detection, image and video analysis, face tracking, Grad-CAM, timelines, evidence frames, downloadable PDF forensic report, history and deletion, deployed as a static frontend plus a hosted API.
\item[Experimental contribution.] A reproducible Kaggle training and evaluation workflow with pre-registered protocols; four models compared under identical code (Table~\ref{tab:fourvideo}); experiments E1--E4 including a frozen promotion rule; E4 as the latest model, with its limits reported (Section~\ref{sec:e4}).
\end{description}

\subsection{What remains limited}
No new detector architecture, no validation of the video aggregation rule on labelled videos, no calibration, no claim about AI-generated video, E4 not approved by the promotion rule, and the live deployment serves v2 (Chapter~\ref{ch:results}, Section~\ref{sec:limits}).


